\chapter{Introduction}\label{ch:introduction}

An interactive vector map runs two paths: a per-tile path that fetches, decodes, and binds the style to each tile's features at load time, and a per-frame path that re-evaluates view-dependent expressions and submits draw calls to the \ac{GPU}.
A slow network stalls the first.
A battery- or thermally-bound device stalls the second.
MapLibre GL~JS\footnote{\url{https://maplibre.org/maplibre-gl-js/docs/}, last accessed 2026-05-28} and Mapbox GL~JS\footnote{\url{https://www.mapbox.com/mapbox-gljs}, last accessed 2026-05-28} deliver these maps by streaming tiled geometry to the browser.
They empower a client with a client-side style document driving a \ac{WebGL} or \ac{WebGPU} based rendering pipeline.
This architecture is what makes panning, zooming, and rotation smooth when it works, and what makes them stutter when it does not.

As map deployments scale from personal projects to platforms serving millions of daily users, the cost of suboptimal performance compounds.
Larger tiles increase transfer time and bandwidth cost, redundant style layers waste rendering cycles and unnecessary properties inflate decoding overhead.
On mobile devices, these inefficiencies translate directly into battery drain and thermal throttling, though we leave direct energy measurement outside the scope of this work.
On unstable network connections, such as on trains or aeroplanes, these translate to degraded user experiences.

Despite this, the map rendering and geospatial community has largely treated style authoring and data preparation as independent concerns: styles are hand-tuned for visual appearance, while tiles are encoded with general-purpose content and compression.
The interaction between the two is left unexploited.
For example, a style filter that matches on a property value may, after constant folding with tile-level statistics, simplify to a form that no longer references this property.
This in turn allows the data encoder to drop the corresponding column from every tile, a saving that neither style simplification nor data pruning could achieve independently.
We ask whether such cross-level interactions produce measurably super-additive gains, or whether the levels compose merely additively, as a central empirical question of this work that we formalise in \cref{sec:problem_statement}.

\section{Research Goals}\label{sec:research_goals}

We investigate automatic, style- and data-driven co-optimization for MapLibre-based map rendering pipelines.
Three research questions guide our work:

\begin{enumerate}[label=\textbf{R\arabic*}]
    \item To what extent \textbf{can automatic data and style co-optimization improve rendering performance and reduce map tile size}?
    \item \textbf{Which techniques} (such as expression simplification, dead layer elimination, layer merging, tile shaving, and columnar re-encoding) \textbf{are most effective}, \textbf{and how do they interact}?  In particular, R2 asks whether the composition of style-level and data-level optimizations is additive or synergistic.
    \item \textbf{How expensive are these optimizations} in terms of preprocessing time, and what are the \textbf{tradeoffs between decoding latency, rendering throughput, and transfer size}?
\end{enumerate}

To answer these questions, we design and implement a standalone optimizer.
It consumes a MapLibre style document and optional tile-level statistics, and emits a replacement $(S', T')$ that is visually equivalent (\cref{eq:visual_equivalence}) and strictly idempotent (\cref{eq:idempotency}).
The magnitude of the resulting gains varies by style and scenario, as quantified in \cref{ch:experiments}.

\newpage
\section{Contributions}\label{sec:contributions}

We make the following main contributions:

\begin{enumerate}[label=\textbf{C\arabic*}]
    \item A \textbf{formalization of joint style and tile co-optimization} as a transformation $\mathcal{O}: \mathcal{S} \times \mathcal{T} \to \mathcal{S} \times \mathcal{T}$ under visual equivalence (\cref{eq:visual_equivalence}) and strict idempotency (\cref{eq:idempotency}).
    Compiler optimization and columnar encoding have each addressed one side in isolation.
    No prior work, to our knowledge, treats both sides jointly under a soundness criterion (\cref{sec:problem_statement}).

    \item A \textbf{novel three-level co-optimization strategy} (\cref{sub:pipeline}).
    The first level applies expression-level peephole rewrites and iterates them to a fixpoint.
    The second level performs structural transformations over a typed style model (dead layer elimination, metadata refinement, layer merging).
    The third level transforms tile payloads and introduces two techniques that are, to our knowledge, novel.
    \textbf{Style-driven tile shaving} (\cref{sub:tile_shaving}) derives a pruning advisory from the optimized style and removes unused properties, geometry types, values, and feature IDs from tile data.
    \textbf{Adaptive multi-strategy encoding} for the \ac{MLT} columnar format (\cref{sub:data_optimizations}) competes sort orderings, integer encodings, and string compression strategies at encode time rather than committing to a fixed plan.
    Style-level passes further refine the pruning advisory, though the cross-family interaction stays additive rather than synergistic (\cref{sub:mlt_discussion}).

    \item An \textbf{end-to-end implementation} of the strategy in Rust and Java.
    It integrates the MapLibre style runtime and the \ac{MLT} tile pipeline, and is released as open source.

    \item An \textbf{empirical evaluation on four complementary corpora}.
    The Germany OpenMapTiles tileset provides an encoder-only, style-independent baseline.
    The ten-style OpenMapTiles deep corpus carries the end-to-end browser benchmarks.
    An eleven-style \ac{OMT}-compatible subset is measured for tile-shaving effectiveness only.
    A thirteen-style Maputnik generalization corpus is evaluated under static style passes alone.
    Together they span 18 scenarios and 19 cumulative ablation steps under fixed pass ordering (interaction effects discussed in \cref{sec:threats}).
    We quantify load-time, tile-volume, and \ac{FPS} effects, and answer the additive-versus-synergistic question of R2 at both the aggregate and per-scenario levels.
\end{enumerate}

\section{Outline}\label{sec:outline}

The pipeline composes expression-level rewrites, structural style passes, and a style-driven pruning advisory that feeds \ac{MLT} re-encoding. \Cref{ch:methods} states the joint problem formally under the constraint of \cref{eq:visual_equivalence,eq:idempotency} and details the three levels, with the prior-art positioning and theoretical background in \cref{ch:related_work} and \cref{ch:theoretical_background}.
\Cref{ch:experiments} measures the isolated and combined effects through a cumulative ablation across 15~styles and 18~scenarios, and \cref{ch:conclusion} revisits R1-R3 in \cref{sec:rq_revisited} and names tile shaving and columnar re-encoding as the primary levers.
